\section{Induction on the scale: reducing pinned averages to smaller instances}
\label{sec:scale-comparison}

\paragraph{What this section proves and how to read it.}
This section proves that, at a least counterexample, every pinned
average (an average with labels fixed at a few atoms and inputs
restricted to residue classes, called a \emph{structured energy})
is at most its price times $Y$, and that the induction then converts
$\Theta<1$ into $s(N)\le CN^{1-a/71}$ in the baseline
($\theta=1$). The mechanism is the single rescaling already used in
\cref{sec:mechanism}: fixing labels restricts $A$ to shorter
square-step progressions, where the induction hypothesis holds.
Inputs: $Y=\Yf(\Lsc,A)$ of \cref{def:record:functional} with atom
laws satisfying (R1)--(R5). Outputs: the pinned-average bound, the
rules (\emph{admissible for the energy estimates},
\cref{def:admissible}---one of three qualified uses of
``admissible'' in the paper) for applying sharp bounds without
counting a price twice, negligibility of forced pairs and of large
pinnings, and \cref{thm:record:lineardensity}. Part~III uses only
these outputs. A first reading should follow the rescaling in
\cref{subsec:scale:fibres} and the induction in
\cref{subsec:scale:induction}; the folded-energy order of operations
can be read as a specified order (sharp bound, then average over
shared labels, then rescale).

\paragraph{Remainders linear in the density.}
Two features of the estimates matter. First, every analytic remainder is
$O(\Lsc^{-\sigma}\alpha)$, with the density $\alpha=|A|/N$ to the first
power. The density enters through the moment bounds of
\cref{lem:record:lift-moments}, which are linear in $\alpha$: an $L^{2r}$
norm of an input costs $\alpha^{1/(2r)}$, and the density exponents of every
H\"older split used below add up to exactly one ($1/2+71/142$ in the
comparison, $2\cdot2/4$ and $70\cdot2/140$ in the cancellation, $72/72$ for
large products). At a counterexample, where $Y>C_0\Lsc^{-a}\alpha^\theta$, a
remainder $C\Lsc^{-\sigma}\alpha$ is at most
$(C/C_0)\Lsc^{a-\sigma}\alpha^{1-\theta}Y$, uniformly in $\alpha$, and the
induction closes at the density exponent $\theta=1$ under the single
condition $a<\sigma$. An absolute remainder $C\Lsc^{-\sigma}$ would be small compared
with $Y$ only if $\alpha$ were not too small, which forces a sparseness
condition and $\theta<1$ (\cref{rem:record:theta-ladder}).

Second, the averaging is done in a fixed order. A sharp operator bound is
applied at fixed shared labels, with a constant that does not depend on
them; the Cauchy--Schwarz inequality is then integrated over the law of the
shared labels; and only then is the comparison applied, to the two resulting
unconditional energies (\cref{lem:record:folded-energy}). No comparison is
pointwise in the selected labels, and no conditional
distribution of the non-shared labels of one half, given the
shared labels, is assumed to be uniform.

The induction hypothesis is used only through the rescaling of
\cref{lem:record:fibre-scale}. Fixing the labels at a set of atoms fixes a
residue class; splitting the class into progressions with a square step
returns the same problem at a shorter scale, because a square-difference-free
set meets a progression of square step in a rescaled square-difference-free
set.

\subsection{Numerical exponents and the inequalities they must satisfy}\label{subsec:scale:stage}

The exponents of the analytic estimates are fixed once.

\begin{definition}[Analytic stage]\label{def:record:analytic-stage}
$\beta=1/20$ is the exponent of the functional; $\eta=17/100$ is the
major-arc exponent of \cref{app:kernel20}; $\gamma=255/1024$ is the
one-slot and pair mixing exponent at the tail primes $p\ge\Pfive$; selected
sets of field primes have product at most $N^\tau$, $\tau=1/201$;
$T_0=N^{t_{\rm an}}$, $t_{\rm an}=1/60$, is the truncation of pair supports
in the cancellation; $H=N^{\beta/2}=N^{1/40}$ is the outside cutoff of the
comparison and of the kernel; $\sigma=1/825$ is the error exponent;
$\delta=10^{-10}$ is the moment loss per slot and $\ell_0=10^{-8}$ the
moment loss of the large-product estimate. A \emph{constant} depends only on
these parameters, on $\Kint$, $\Mzero$, $Q_{\rm pre}$, the reference laws and
the moment constants of \cref{lem:record:lift-moments}; it never depends on
$\Lsc$, $A$, the selected atoms, their labels, the intervals or the residues.
\end{definition}

\paragraph{The stage inequalities.}
\Cref{cert:record:analytic-stage} (\cref{app:certificates}) checks in
exact rational arithmetic that these parameters satisfy the following strict
inequalities:
\begin{enumerate}[label=(S\arabic*)]
\item $\gamma\beta/2-\tau-72\delta-\sigma>0$;
\item $\gamma t_{\rm an}-72\delta-\sigma>0$;
\item $\beta/6-t_{\rm an}(1/2-\gamma)-72\delta-\sigma>0$;
\item $\tau/4-\ell_0-\sigma>0$;
\item $1-2\tau-6\beta>0$, $\beta/2-\tau>0$, $\beta/2-t_{\rm an}>0$ and
 $\beta(1-2\tau)-\beta/2>0$;
\item the seven kernel slacks of \cref{lem:record:arcs}, computed in
 \cref{cor:kernel20}, are positive;
\item the tail coefficients and the tail gates of
 \cref{lem:record:tail-mixing} hold; only the gate at $\Pfive$ is used.
\end{enumerate}
Here (S1) is used in the comparison, where only the weaker
$\gamma\beta/2-72\delta>\sigma$ is needed for an individual averaged estimate
(\cref{lem:record:density-dc}); its full slack also absorbs the sum of raw
comparison errors in \cref{prop:density:production}. Conditions (S2) and (S3) are used in the accepted-pair
cancellation (\cref{lem:record:density-can}), where they are the two saving
exponents $\gamma t_{\rm an}=17/4096$ and
$1/120-t_{\rm an}(1/2-\gamma)=17/4096$ less the moment loss; (S4) in the
large-product estimate (\cref{lem:record:density-lp}); the four inequalities
of (S5), in this order, give the moment hypothesis at the target scale,
complete selected fibres ($HD_E\le Q$ when $q_E\le N^\tau$), complete pair
fibres and the kernel modulus ($HN^{t_{\rm an}}Q_{\rm pre}\le Q$), and the
target cutoff ($n^\beta\ge H$), through \cref{lem:scale-conditions}; (S6) is
the admissibility of the kernel triple of \cref{cor:kernel20}, used in
\cref{app:kernel20}, and (S7) is used in \cref{lem:record:tail-mixing}. None
of (S1)--(S7) involves $\Kint$ or the clique prefix. Apart from the slot
prices $\xi_v$ and the density factors $\kappa_v$ of
\cref{def:record:structured-energy}, in which $\Kint$, divisors of $\Mzero$
and the clique powers $p^{e_p}$ appear explicitly, the fixed integers
$\Kint$, $\Mzero$, $Q_{\rm cl}$ and $Q_{\rm pre}$ enter the estimates of this
section only through constants, through the term $u_\Lsc\to0$ and through the
starting scale of \cref{lem:scale-conditions}; this is possible because every
inequality is strict.

\subsection{Residue-class restriction and square-step rescaling}\label{subsec:scale:fibres}

For a set $E$ of atoms of $\mathcal A_N$ we write $X_E$ for the product of
their groups, $D_E=|X_E|$ for its full order, $q_E$ for the product of the
primes of the atoms of $E$ with $e_p=1$, and $E_{\rm cl}=E\cap\mathcal P_1$
for its clique primes. Thus $D_E=q_E\prod_{p\in E_{\rm cl}}p^{e_p}\le
q_EQ_{\rm cl}$. For a divisor $D_0$ of $\Mzero$ put $D=D_0D_E$. The two
factors are coprime: the primes of $\Mzero$ are $2$ and the primes
$p\equiv3\pmod4$ below $Q_3$, none of which is a prime of an atom. The
\emph{square step} of $E$ and $D_0$ is the least perfect square divisible by
$D$ \eqref{eq:sq},
\begin{equation}\label{eq:square-step}
 R=\sq(D)=\sq(D_0)\,q_E^2\prod_{p\in E_{\rm cl}}p^{e_p};
\end{equation}
the second form holds because $\sq$ is multiplicative on coprime factors,
$q_E$ is squarefree and every $e_p$ is even. The primes of $R$ are those of
$D$, so $R$ is coprime to every prime of an atom outside $E$.

\begin{lemma}[Exact fibres and the square-step comparison input]
\label{lem:record:fibre-scale}
Let $\Lsc\ge2$, $N=\Kint\Lsc$, let $A\subseteq[N]$ be square-difference-free
and let $E\subseteq\mathcal A_N$, with $\mathcal S=\mathcal A_N\setminus E$.
Let $I\in\mathcal I_\Lsc$, $D_0\mid\Mzero$, $s\in\Z/D_0\Z$, $z\in X_E$ and
$g=g^{I,D_0,s}$. Put $D=D_0D_E$, let $\mathfrak a\bmod D$ be the class with
$\mathfrak a\equiv s\pmod{D_0}$ and $\mathfrak a\equiv z_p\pmod{p^{e_p}}$ for the primes $p$ of
$E$, and let
\[
 \alpha_{\rm f}=\frac D{|I|}\,\#\{x\in A\cap I:\ x\equiv\mathfrak a\ (\mathrm{mod}\ D)\}
\]
be the density of $A$ in the fibre of $z$.
\begin{enumerate}[label=(\alph*)]
\item \emph{(Exact fibre.)} For every $f:I\to\C$ and every character $\xi$ of
 $X_{\mathcal S}$ with $q(\xi)D_E\le Q$, the coefficient at $\xi$ of the
 specialization $(\mathcal L_{I,N}f)(z,\cdot)$ is
 $(D_E/|I|)\sum_{x\in I,\,x\equiv z\,(D_E)}f(x)e(-x\xi)$, where
 $x\equiv z\ (D_E)$ means $x\equiv z_p\pmod{p^{e_p}}$ for each prime $p$ of
 $E$. In particular
 \[
  \widehat{g(z,\cdot)}(\xi)=\frac D{|I|}\sum_{\substack{x\in A\cap I\\ x\equiv\mathfrak a\,(D)}}e(-x\xi).
 \]
\item \emph{(Square step.)} Let $R=\sq(D)$ be the square step of $E$ and
 $D_0$, $\ell=\lceil|I|/(\Kint R)\rceil$, $n=\Kint\ell$ and
 $\lambda=Rn/|I|$. Write $I=c_I+[|I|]$, $A_I=(A\cap I)-c_I\subseteq[|I|]$ and
 $J_{\mathfrak a}=\{j\in[R]:j\equiv\mathfrak a-c_I\pmod D\}$ (so $|J_{\mathfrak a}|=R/D$), and for $j\in J_{\mathfrak a}$
 put $c_j=j-R$,
 \[
  A'_j=\{n'\in[n]:\ Rn'+c_j\in A_I\},\qquad
  T_j(y)=R^{-1}(y-c_I-c_j)\quad(y\in X_{\mathcal S}),
 \]
 $h_j=\mathcal L^{\mathcal S}_{[n],n^\beta}\one_{A'_j}$, the lift of
 $\one_{A'_j}$ on $X_{\mathcal S}$ with cutoff $n^\beta$, and
 \[
  \widetilde g=\frac{\lambda D}{R}\sum_{j\in J_{\mathfrak a}}h_j\circ T_j ,
 \]
 the \emph{square-step model} of $g(z,\cdot)$. Each $A'_j$ is
 square-difference-free, each $T_j$ is a map of the form in
 \ref{it:record:R3}, and
 \[
  \frac DR\sum_{j\in J_{\mathfrak a}}\alpha_j=\frac{\alpha_{\rm f}}\lambda,\qquad
  \alpha_j=\frac{|A'_j|}n .
 \]
 If $H\le n^\beta$ and $HD_E\le Q$, then $\widetilde g$ and $g(z,\cdot)$
 have the same Fourier coefficient at every character $\xi$ of
 $X_{\mathcal S}$ with $q(\xi)\le H$.
\item \emph{(Restoration.)} $Z_{\mathcal S}(h_j)\le\Yf(\ell,A'_j)$ for every
 $j\in J_{\mathfrak a}$.
\item \emph{(Scales.)} $1\le\lambda<1+\Kint R/|I|$ and $\ell\ge\Lsc/\xi_I$
 with $\xi_I=(N/|I|)R$. Moreover $\ell<\Lsc$, unless $E=\emptyset$,
 $D_0=1$ and $I=[N]$.
\item \emph{(Moments.)} If $n\ge10n^{6\beta}$, then for every integer
 $r\ge1$ and every $\epsilon>0$,
 $\|\widetilde g\|_{2r}^{2r}\le C_{r,\epsilon}N^\epsilon\lambda^{2r-1}\alpha_{\rm f}$,
 with the constant $C_{r,\epsilon}$ of \cref{lem:record:lift-moments}(c).
\end{enumerate}
\end{lemma}

\begin{proof}
(a) Every prime of $D_0$ divides $\Mzero$ and none divides $D_E$, so the
class $\mathfrak a$ exists and is unique. Write a character of $G_N$ as $\xi+\zeta$
with $\xi$ a character of $X_{\mathcal S}$ and $\zeta$ one of $X_E$; their
conductors are coprime, so $q(\xi+\zeta)=q(\xi)q(\zeta)$, and $q(\zeta)$
divides $D_E$. If $q(\xi)D_E\le Q$, every character $\zeta$ of $X_E$ occurs
together with $\xi$ in the lift, so the coefficient of $e(y\cdot\xi)$ in
$(\mathcal L_{I,N}f)(z,\cdot)$ is
\[
 \frac1{|I|}\sum_{x\in I}f(x)e(-x\xi)\sum_{\zeta}e((z-x)\cdot\zeta)
 =\frac{D_E}{|I|}\sum_{\substack{x\in I\\ x\equiv z\,(D_E)}}f(x)e(-x\xi),
\]
by orthogonality of the characters of each cyclic group $\Z/p^{e_p}\Z$. For
$f=D_0\one_{x\equiv s\,(D_0)}\one_A$ the Chinese remainder theorem
gives the second formula, and it holds also when the fibre is empty.
Completeness of the selected group is necessary here: for the cyclic group
$\Z/9\Z$ with cutoff $3$, $I=[1,9]$, $A=\{9\}$ and label $3$, the retained
characters have numerators $0,3,6$ and the specialized lift equals $1/3$,
although the fibre $\{x\equiv3\ (9)\}\cap A$ is empty.

(b) Since $D$ divides $R$, exactly $R/D$ residues $j\in[R]$ satisfy
$j\equiv\mathfrak a-c_I\pmod D$. The map $(n',j)\mapsto Rn'+c_j=R(n'-1)+j$ is a
bijection $[n]\times[R]\to[Rn]$, and $[|I|]\subseteq[Rn]$ because
$Rn\ge|I|$; restricted to $j\in J_{\mathfrak a}$ it is a bijection onto the integers of
$[Rn]$ in the class $\mathfrak a-c_I\bmod D$. The integers of $(|I|,Rn]$ do not lie in
$A_I$, so the padding contributes nothing, and
$\sum_{j\in J_{\mathfrak a}}|A'_j|=\#\{x\in A\cap I:x\equiv\mathfrak a\ (D)\}$. Therefore
$(D/R)\sum_j\alpha_j=D\,\#\{x\in A\cap I:x\equiv\mathfrak a\ (D)\}/(Rn)
=\alpha_{\rm f}|I|/(Rn)=\alpha_{\rm f}/\lambda$. If $n'_1<n'_2$ lie in
$A'_j$ and $n'_2-n'_1=k^2$ with $k\ge1$, then two elements of $A$ differ by
$Rk^2=(\sqrt R\,k)^2$, which is impossible because $R$ is a perfect square.
The integer $R$ is coprime to every prime of $\mathcal S$, and as $R$ is a
perfect square, $R^{-1}$ is a square unit modulo every $p^{e_p}$, $p$ a prime
of $\mathcal S$; so $T_j$ is of the form in \ref{it:record:R3}. Let $\xi$ be
a character of $X_{\mathcal S}$ with $q(\xi)\le H$. Multiplication by the
unit $R$ permutes the characters of $X_{\mathcal S}$ and preserves
conductors, and $e(T_j(y)\cdot R\xi)=e((y-c_I-c_j)\cdot\xi)$. Since
$q(R\xi)=q(\xi)\le H\le n^\beta$, the character $R\xi$ is retained in $h_j$,
and the coefficient of $e(y\cdot\xi)$ in $\widetilde g$ is
\[
 \frac{\lambda D}{R}\sum_{j\in J_{\mathfrak a}}e(-(c_I+c_j)\xi)\,
 \frac1n\sum_{n'\in A'_j}e(-n'R\xi)
 =\frac{\lambda D}{Rn}\sum_{\substack{x\in A\cap I\\ x\equiv\mathfrak a\,(D)}}e(-x\xi)
 =\frac D{|I|}\sum_{\substack{x\in A\cap I\\ x\equiv\mathfrak a\,(D)}}e(-x\xi).
\]
By (a), which applies since $q(\xi)D_E\le HD_E\le Q$, this is the
coefficient of $g(z,\cdot)$.

(c) Let $\mathcal A_n$ be the atoms of the functional at the scale $\ell$,
whose group has cutoff $n^\beta$ ($\Kint\ell=n$); so
$\mathcal A_n\subseteq\mathcal A_N$, and put $E'=E\cap\mathcal A_n$. A
character with $q\le n^\beta$ involves only primes at most $n^\beta$, and
every such prime of $\mathcal S$ lies in an atom of $\mathcal A_n$ (the
atoms of $\mathcal P_1\cup\mathcal P_3$ belong to every scale), so $h_j$
depends only on the
coordinates of $\mathcal S'=\mathcal S\cap\mathcal A_n=\mathcal A_n\setminus E'$,
and $Z_{\mathcal S}(h_j)=Z_{\mathcal S'}(h_j)$ by Fubini
(\cref{lem:record:holder}(d)). Let $u_j=\mathcal L_{[n],n}\one_{A'_j}$ be the
full lift at scale $n$. Averaging $u_j$ over the coordinates of $E'$ deletes
exactly the characters that are nonconstant on an atom of $E'$ and leaves the
others; the remaining characters are those of $h_j$, with the same
coefficients. Hence $\E_{E'}u_j=h_j$. Averaging over atoms does not increase
$Z$ (\cref{lem:record:holder}(d)), and the atom laws at scale $\ell$ are those
at scale $\Lsc$ by \ref{it:record:R1}, so
\[
 Z_{\mathcal S'}(h_j)=Z_{\mathcal A_n}(\E_{E'}u_j)\le Z_{\mathcal A_n}(u_j)
 =\Yf(\ell,A'_j).
\]

(d) $n=\Kint\lceil|I|/(\Kint R)\rceil$ lies in $[|I|/R,|I|/R+\Kint)$, which
gives the bounds for $\lambda$, and $\ell\ge|I|/(\Kint R)=\Lsc/\xi_I$. If
$D>1$, then $R\ge4$, because $R$ is a perfect square, and
$\ell\le\lceil N/(4\Kint)\rceil=\lceil\Lsc/4\rceil<\Lsc$. If $D=1$, then
$E=\emptyset$, $D_0=1$ and $R=1$; for $I=I_b$,
$\ell=\lceil\Lsc/\Kint\rceil<\Lsc$, since $\Kint\ge2$ and $\Lsc\ge2$.

(e) $T_j$ preserves the uniform measure, so by Minkowski
$\|\widetilde g\|_{2r}\le\lambda\frac DR\sum_j\|h_j\|_{2r}$. Each $h_j$ is a
whole-support piece of the lift of $\one_{A'_j}$ on $X_{\mathcal S}$ with
cutoff $n^\beta$, and $|[n]|\ge10(n^\beta)^6$, so
$\|h_j\|_{2r}^{2r}\le C_{r,\epsilon}N^\epsilon\alpha_j$ by
\cref{lem:record:lift-moments}(c) and $n\le N$. Concavity of
$x\mapsto x^{1/(2r)}$ and (b) give
$\|\widetilde g\|_{2r}\le\lambda(C_{r,\epsilon}N^\epsilon\alpha_{\rm f}/\lambda)^{1/(2r)}$.
\end{proof}

The hypotheses of (b) and (e) hold for every selection used below, from a
fixed starting scale on. We collect once all the conditions on the scale
that the comparison and the cancellation need.

\begin{lemma}[Scale conditions]\label{lem:scale-conditions}
There is $\Lsc_*\ge2$, depending only on the analytic stage, $\Kint$,
$\Mzero$, $Q_{\rm cl}$ and $Q_{\rm pre}$, such that for every
$\Lsc\ge\Lsc_*$, with $N=\Kint\Lsc$, $Q=N^\beta$, $H=N^{\beta/2}$ and
$T_0=N^{t_{\rm an}}$:
\begin{enumerate}[label=(\roman*)]
\item $\Lsc\ge10N^{6\beta}$;
\item $HQ_{\rm cl}N^\tau\le Q$;
\item $n_0:=N^{1-2\tau}/(\Kint\Mzero^2Q_{\rm cl})$ satisfies
 $n_0^\beta\ge H$ and $n_0\ge10N^{6\beta}$;
\item $\Kint^2\Mzero^2Q_{\rm cl}N^{2\tau-1}\le1$;
\item $N^\beta\ge\Pfive$, so that every prime of $[\Pfour,\Pfive)$ is a
 coordinate of $G_N$;
\item $HQ_{\rm pre}T_0\le Q$ and
 $\Mzero Q_{\rm pre}T_0\le\min(N^{1/20},\Lsc)$;
\item $N$ is at least the thresholds $N_0$ and $N_0'$ of
 \cref{prop:record:kernel20}(i)--(iii) at the triple of
 \cref{cor:kernel20}, $(\beta,\kappa,\eta)=(1/20,1/20,17/100)$, with
 $C_D=C_H=1$, $\mathcal P_0=\mathcal P_1$ and $m_p=e_p/2$.
\end{enumerate}
Consequently, let $E\subseteq\mathcal A_N$ with $q_E\le N^\tau$, $D_0\mid\Mzero$
and $I\in\mathcal I_\Lsc$, and let $R$, $n$ and $\lambda$ be as in
\cref{lem:record:fibre-scale}(b). Then $R\le\Mzero^2Q_{\rm cl}N^{2\tau}$,
$HD_E\le Q$, $n_0\le n\le N$, $n^\beta\ge H$ and $n\ge10n^{6\beta}$; and
$1\le\lambda\le2$ and $\lambda^{72}\le1+u_\Lsc$, where
\[
 u_\Lsc=\big(1+\Kint^2\Mzero^2Q_{\rm cl}N^{2\tau-1}\big)^{72}-1
\]
tends to zero as $\Lsc\to\infty$. In particular the hypotheses of
\cref{lem:record:fibre-scale}(b) and (e) hold.
\end{lemma}

\begin{proof}
Each of (i)--(iv) and (vi) compares two powers of $N$, up to fixed
constants, with a positive gap in the exponents: $1-6\beta>0$ for (i),
$\beta/2-\tau>0$ for (ii), $\beta(1-2\tau)-\beta/2>0$ and $1-2\tau-6\beta>0$
for (iii), $1-2\tau>0$ for (iv), and $\beta/2-t_{\rm an}>0$, hence also
$\beta-t_{\rm an}>0$ and $1-t_{\rm an}>0$, for (vi); all of them follow from
(S5). Conditions (v) and (vii) are fixed thresholds. So all seven hold for
large $\Lsc$.

For the consequences, $\sq(D_0)\le D_0^2\le\Mzero^2$ and
$\prod_{p\in E_{\rm cl}}p^{e_p}\le Q_{\rm cl}$ give
$R\le\Mzero^2Q_{\rm cl}N^{2\tau}$, and $D_E\le q_EQ_{\rm cl}$ with (ii) gives
$HD_E\le Q$. Next $n=\Kint\lceil|I|/(\Kint R)\rceil\le\Kint\lceil N/\Kint\rceil=N$,
and $n\ge|I|/R\ge(N/\Kint)/(\Mzero^2Q_{\rm cl}N^{2\tau})=n_0$ because
$|I|\ge\Lsc=N/\Kint$; by (iii), $n^\beta\ge n_0^\beta\ge H$ and
$n\ge n_0\ge10N^{6\beta}\ge10n^{6\beta}$. Finally
\cref{lem:record:fibre-scale}(d) gives
$1\le\lambda<1+\Kint R/|I|\le1+\Kint^2\Mzero^2Q_{\rm cl}N^{2\tau-1}$, which is
at most $2$ by (iv), and raising to the power $72$ gives
$\lambda^{72}\le1+u_\Lsc$.
\end{proof}

Conditions (i)--(iv) are used in the comparison
(\cref{lem:record:density-dc,prop:record:density-structured}), (i) and
(v)--(vii) in the cancellation (\cref{lem:record:density-can}), and (i) in the
large-product estimate (\cref{lem:record:density-lp}).

\subsection{Pinned averages at a least counterexample}
\label{subsec:scale:structured}

All energy estimates are made at a least counterexample to the bound that
the induction proves.

\begin{definition}[Least normalized counterexample]
\label{def:record:counterexample}
Let $a>0$, $0<\theta\le1$ and $C_0\ge1$. A pair $(\Lsc,A)$, with
$A\subseteq[\Kint\Lsc]$ square-difference-free, is a \emph{least normalized
counterexample} for $(a,\theta,C_0)$ if $\Yf(\Lsc,A)>C_0\Lsc^{-a}\alpha^\theta$,
and $\Yf(\ell,A')\le C_0\ell^{-a}\alpha_{A'}^\theta$ for every $\ell<\Lsc$ and
every square-difference-free $A'\subseteq[\Kint\ell]$, where
$\alpha_{A'}=|A'|/(\Kint\ell)$ and $0^\theta=0$. Then $\alpha>0$, since
$F_\emptyset=0$.
\end{definition}

\begin{definition}[Structured energies and prices]
\label{def:record:structured-energy}
Let $a>0$ and $0<\theta\le1$ be the exponents of the induction. Let
$\Lsc\ge1$, $N=\Kint\Lsc$, let $A\subseteq[N]$ be square-difference-free and
$\alpha=|A|/N>0$. A \emph{structured energy} at $(\Lsc,A)$ is given by the
following data.
\begin{enumerate}[label=(\alph*)]
\item A set $E\subseteq\mathcal A_N$ of \emph{selected} atoms with
 $q_E\le N^\tau$; put $\mathcal S=\mathcal A_N\setminus E$.
\item A finite signed measure $\rho$ on $X_E^V$ whose variation $|\rho|$ has
 every one-slot marginal equal to $\|\rho\|\Unif(X_E)$, where
 $\|\rho\|=|\rho|(X_E^V)$.
\item For each $v\in V$ a divisor $D_{0,v}$ of $\Mzero$ and a length
 $M_v\in\{\Lsc,N\}$, and a probability law $\mathbf P$ on index vectors
 $\omega=(I_v,s_v)_{v\in V}$ with $I_v\in\mathcal I_\Lsc$, $|I_v|=M_v$ and
 $s_v\in\Z/D_{0,v}\Z$.
\end{enumerate}
The energy is
\[
 \mathfrak E=\E_{\omega\sim\mathbf P}\int\Lambda_{\mu_{\mathcal S}}
 \big((g_v^\omega(z_v,\cdot))_{v\in V}\big)\,d\rho(z),\qquad
 g_v^\omega=g^{I_v,D_{0,v},s_v}.
\]
A slot is \emph{unchanged} if $E=\emptyset$, $D_{0,v}=1$ and $M_v=N$ (then
$g_v^\omega=F_A$), and \emph{changed} otherwise. With
$R_E=q_E^2\prod_{p\in E_{\rm cl}}p^{e_p}$, the slot prices are
\[
 \xi_v=\frac N{M_v}\,\sq(D_{0,v})R_E,\qquad
 \bar\alpha_v=\E_{\mathbf P}\,\frac{D_{0,v}}{M_v}\#\{x\in A\cap I_v:
 x\equiv s_v\ (D_{0,v})\}.
\]
The \emph{main price} is
$\Pi_{\rm main}=\prod_v\xi_v^{a/72}\prod_{v\ \rm changed}(\bar\alpha_v/\alpha)^{\theta/72}$.
If $\kappa_v\ge1$ are numbers with $\bar\alpha_v\le\kappa_v\alpha$ for every
changed slot, the \emph{coarse price} is
$\Pi=\prod_v\xi_v^{a/72}\prod_{v\ \rm changed}\kappa_v^{\theta/72}$; thus
$\Pi\ge\max(1,\Pi_{\rm main})$.
\end{definition}

One may take $\kappa_v=1$ if $(I_v,s_v)$ is uniform on
$\{I_1,\dots,I_\Kint\}\times\Z/D_{0,v}\Z$, or if $M_v=N$ and $s_v$ is
uniform; $\kappa_v=\Kint$ if $s_v$ is uniform given a fixed interval $I_v$,
since then $\bar\alpha_v=\alpha_{I_v}\le\Kint\alpha$; and always
$\kappa_v=\Mzero\Kint$. Since $\sq(D_{0,v})R_E=\sq(D_{0,v}D_E)$, the price
$\xi_v$ is $\xi_{I_v}$ of \cref{lem:record:fibre-scale}(d): for $M_v=N$ it is
the square step of $E$ and $D_{0,v}$, and for $M_v=\Lsc$ it is $\Kint$ times
it. If a clique prime $p$ belongs to $E$, its contribution to
$\prod_v\xi_v^{a/72}$ over the $72$ slots is $p^{e_pa}$; a prefix
atom not in $E$ contributes no factor.

\begin{lemma}[Averaged comparison with a density factor]
\label{lem:record:density-dc}
There is a constant $C_{\rm dc}$ with the following property. Let
$\Lsc\ge\Lsc_*$ (\cref{lem:scale-conditions}) and consider a structured
energy at $(\Lsc,A)$. Put $H=N^{1/40}$. For every $\omega$, $z\in X_E^V$ and
changed slot $v$, let $\widetilde g_v^\omega(z_v,\cdot)$ be the square-step
model of \cref{lem:record:fibre-scale}(b) for
$(E,I_v,D_{0,v},s_v,z_v)$; for an unchanged slot put
$\widetilde g_v^\omega=F_A$. Let $\widetilde{\mathfrak E}$ be the energy with
$\widetilde g_v^\omega$ in place of $g_v^\omega$. Then
\begin{multline*}
 \E_{\mathbf P}\int\big|\Lambda_{\mu_{\mathcal S}}((g_v^\omega(z_v,\cdot))_v)
 -\Lambda_{\mu_{\mathcal S}}((\widetilde g_v^\omega(z_v,\cdot))_v)\big|\,d|\rho|(z)\\
 \le C_{\rm dc}\|\rho\|N^{-\gamma\beta/2+72\delta}\alpha
 \le C_{\rm dc}\|\rho\|\Lsc^{-\sigma}\alpha ;
\end{multline*}
in particular $|\mathfrak E-\widetilde{\mathfrak E}|\le
C_{\rm dc}\|\rho\|\Lsc^{-\sigma}\alpha$.
\end{lemma}

\begin{proof}
\emph{Step 0: the scale.} By \cref{lem:scale-conditions}, applied to $E$ and
each $D_{0,v}$, $I_v$, the hypotheses $H\le n^\beta$, $HD_E\le Q$ and
$n\ge10n^{6\beta}$ of \cref{lem:record:fibre-scale}(b),(e) hold at every
changed slot, and $1\le\lambda_v\le2$.

\emph{Step 1: telescoping.} Fix $\omega$ and $z$, and let $v_1,\dots,v_k$ be
the changed slots. Then
$\Lambda((g_v)_v)-\Lambda((\widetilde g_v)_v)=\sum_{i=1}^k\Lambda(Y^{(i)})$,
where $Y^{(i)}$ carries $\widetilde g_{v_j}$ at $v_j$ for $j<i$,
$g_{v_i}-\widetilde g_{v_i}$ at $v_i$, $g_{v_j}$ at $v_j$ for $j>i$, and $F_A$
at the unchanged slots (all specialized at $z$). By
\cref{lem:record:fibre-scale}(b), the input at $v_i$ has no character of
conductor at most $H$. The law $\mu_{\mathcal S}$ is the product of the laws
of the atoms of $\mathcal S$ below $\Pfive$, a probability measure with
uniform one-slot marginals by \ref{it:record:R3}, and of the tail laws
$\mu_p$, $p\ge\Pfive$, which satisfy the mixing hypothesis
\ref{it:record:R4} at every slot (in the notation of
\cref{def:mixing}). Every character of the first factor has conductor
dividing $Q_{\rm pre}$. By \cref{lem:record:prefix-conductor} with
$C'=Q_{\rm pre}$,
\[
 |\Lambda(Y^{(i)})|\le\Big(\frac{Q_{\rm pre}}H\Big)^{\gamma}
 \big\|(g_{v_i}-\widetilde g_{v_i})(z_{v_i},\cdot)\big\|_2
 \prod_{w\ne v_i}\|Y^{(i)}_w(z_w,\cdot)\|_{142}.
\]

\emph{Step 2: averaging.} Integrate over $\omega\sim\mathbf P$ and
$z\sim|\rho|/\|\rho\|$ and apply H\"older's inequality with the exponents
$2$ and $142$ ($71$ times), so that $1/2+71/142=1$:
\[
 \E\Big[\|X\|_2\prod_{w}\|Y_w\|_{142}\Big]\le(\E\|X\|_2^2)^{1/2}
 \prod_w(\E\|Y_w\|_{142}^{142})^{1/142}.
\]
Under $|\rho|/\|\rho\|$ each $z_w$ is uniform on $X_E$ and independent of
$\omega$. Hence, by Fubini and \cref{lem:record:lift-moments}(c) (the lift on
$G_N$ is a whole-support piece of the kind treated there, with
$\mathcal S'=\mathcal P_1$, and $|I_w|\ge\Lsc\ge10N^{6\beta}$), for
$2r\in\{2,142\}$ and every $\epsilon>0$,
\[
 \E\|g_w^\omega(z_w,\cdot)\|_{2r}^{2r}=\E_\omega\|g_w^\omega\|_{L^{2r}(G_N)}^{2r}
 \le C_{r,\epsilon}N^\epsilon\Mzero^{2r}\,\E_\omega\alpha_{I_w}
 \le C_{r,\epsilon}N^\epsilon\Mzero^{2r}\Kint\alpha ,
\]
because $\alpha_I\le\Kint\alpha$ for $|I|\ge N/\Kint$; the same holds for
$F_A$. By \cref{lem:record:fibre-scale}(e), with $\lambda_w\le2$,
$\E\|\widetilde g_w(z_w,\cdot)\|_{2r}^{2r}\le C_{r,\epsilon}N^\epsilon2^{2r-1}
\E\alpha_{{\rm f},w}$, where $\alpha_{{\rm f},w}$ is the fibre density of
$z_w$, and
$\E_z\alpha_{{\rm f},w}=D_{0,w}\#\{x\in A\cap I_w:x\equiv s_w\}/|I_w|
\le\Mzero\Kint\alpha$, because the class $\mathfrak a(z_w)$ is uniform among the $D_E$
classes modulo $D$ compatible with $s_w$. Finally
$\|X\|_2^2\le2\|g\|_2^2+2\|\widetilde g\|_2^2$. Thus every factor is at most
$(C_\epsilon N^\epsilon\alpha)^{1/(2r)}$, and each of the at most $72$ terms
is bounded by $Q_{\rm pre}^\gamma H^{-\gamma}C_\epsilon N^{\epsilon}\alpha$,
the density exponents adding to $1/2+71/142=1$. Take $\epsilon=72\delta$.
Finally $N^{-\gamma\beta/2+72\delta}\le N^{-\sigma}\le\Lsc^{-\sigma}$ by (S1).
\end{proof}

\begin{proposition}[Density-structured energy estimate]
\label{prop:record:density-structured}
Let $0<\theta\le1$, $a>0$ and $C_0\ge1$, and let $(\Lsc,A)$ be a least
normalized counterexample for $(a,\theta,C_0)$
(\cref{def:record:counterexample}), with $\Lsc\ge\Lsc_*$ and
$Y=\Yf(\Lsc,A)$. Then every structured energy at $(\Lsc,A)$ satisfies
\[
 |\mathfrak E|\le\|\rho\|\big[(1+u_\Lsc)\,\Pi_{\rm main}\,Y
 +C_{\rm dc}\Lsc^{-\sigma}\alpha\big]
 \le\|\rho\|\Big(1+u_\Lsc+\frac{C_{\rm dc}}{C_0}\Lsc^{a-\sigma}
 \alpha^{1-\theta}\Big)\Pi\,Y,
\]
with $u_\Lsc\to0$ as in \cref{lem:scale-conditions}.
\end{proposition}

\begin{proof}
By \cref{lem:record:density-dc} it suffices to bound $\widetilde{\mathfrak E}$
by $\|\rho\|(1+u_\Lsc)\Pi_{\rm main}Y$. Fix $\omega$ and $z$. By the signed
H\"older inequality of \cref{lem:record:holder}(a) for $\mu_{\mathcal S}$,
$|\Lambda_{\mu_{\mathcal S}}((\widetilde g_v)_v)|\le
\prod_vZ_{\mathcal S}(\widetilde g_v(z_v,\cdot))^{1/72}$. At an unchanged slot
$E=\emptyset$, $\mathcal S=\mathcal A_N$ and $Z_{\mathcal S}(F_A)=Y$. At a
changed slot $v$, with the data $D,R,J_{\mathfrak a},h_j,A'_j,\lambda_v,\ell_v$ of
\cref{lem:record:fibre-scale}(b) for that slot, the triangle inequality for
the seminorm $Z_{\mathcal S}^{1/72}$ and its invariance under the maps $T_j$
(\cref{lem:record:holder}(a),(c)), then \cref{lem:record:fibre-scale}(c), give
\[
 Z_{\mathcal S}(\widetilde g_v)^{1/72}\le\lambda_v\frac DR\sum_{j\in J_{\mathfrak a}}
 Z_{\mathcal S}(h_j)^{1/72}\le\lambda_v\frac DR\sum_{j\in J_{\mathfrak a}}
 \Yf(\ell_v,A'_j)^{1/72}.
\]
Here $\ell_v<\Lsc$ by \cref{lem:record:fibre-scale}(d), so the hypothesis on
smaller scales gives $\Yf(\ell_v,A'_j)\le C_0\ell_v^{-a}\alpha_j^\theta$, also
when $A'_j=\emptyset$. Concavity of $x\mapsto x^{\theta/72}$ and
\cref{lem:record:fibre-scale}(b) give
$\frac DR\sum_j\alpha_j^{\theta/72}\le(\alpha_{{\rm f},v}/\lambda_v)^{\theta/72}$,
where $\alpha_{{\rm f},v}$ is the fibre density of $z_v$, hence
\[
 Z_{\mathcal S}(\widetilde g_v)\le\lambda_v^{72-\theta}C_0\ell_v^{-a}\alpha_{{\rm f},v}^\theta
 \le(1+u_\Lsc)C_0\Lsc^{-a}\xi_v^a\alpha_{{\rm f},v}^\theta ,
\]
using $\ell_v\ge\Lsc/\xi_v$ and $\lambda_v^{72}\le1+u_\Lsc$
(\cref{lem:record:fibre-scale}(d), \cref{lem:scale-conditions}). The price
$\xi_v$ does not depend on $\omega$ or $z$. Every factor is nonnegative, so
H\"older's inequality with $72$ equal exponents on the probability space of
$(\omega,z)$, with law $\mathbf P\otimes|\rho|/\|\rho\|$, gives
\[
 |\widetilde{\mathfrak E}|\le\|\rho\|\prod_v\big(\E\,Z_{\mathcal S}
 (\widetilde g_v(z_v,\cdot))\big)^{1/72}.
\]
At a changed slot, Jensen's inequality ($\theta\le1$) and
$\E\alpha_{{\rm f},v}=\bar\alpha_v$ (as in Step~2 of
\cref{lem:record:density-dc}) give $\E Z_{\mathcal S}(\widetilde g_v)\le
(1+u_\Lsc)C_0\Lsc^{-a}\alpha^\theta\xi_v^a(\bar\alpha_v/\alpha)^\theta\le
(1+u_\Lsc)Y\xi_v^a(\bar\alpha_v/\alpha)^\theta$, since
$C_0\Lsc^{-a}\alpha^\theta<Y$. Multiplying over the slots proves the first
inequality. For the second, $\Pi_{\rm main}\le\Pi$, $\Pi\ge1$ and
$C_{\rm dc}\Lsc^{-\sigma}\alpha=(C_{\rm dc}/C_0)\Lsc^{a-\sigma}\alpha^{1-\theta}
\cdot C_0\Lsc^{-a}\alpha^\theta<(C_{\rm dc}/C_0)\Lsc^{a-\sigma}
\alpha^{1-\theta}\,\Pi\,Y$.
\end{proof}

The additive error is bounded before the main term is estimated, and the
main term carries no additive error, so no density ratio is pulled through
an additive error; in particular a density ratio below one is
never factored out of the displayed bound, whose additive error
was estimated before any ratio is formed (an example of the
error this avoids is in \cref{subsec:history:lessons}). The
second form uses only the coarse factors $\kappa_v\ge1$: a fixed interval
index costs at most $\Kint^{\theta/72}$ per changed slot, and the finer
factor $(\alpha_{I_v}/\alpha)^{\theta/72}$ is never multiplied into the
additive term (\cref{rem:wrap-coarse}).

The average over the $R/D$ progressions of the square step is essential, and
their number costs nothing: the weights $D/R$ of the triangle inequality add
up to one, and concavity moves the average inside the power $\theta/72$. A
changed slot therefore costs only $\xi_v^{a/72}$, where $\xi_v$ bounds the
factor by which the child scale $\ell_v\ge\Lsc/\xi_v$ is shorter, together
with the density of its fibre.

\subsection{Differences of laws at one cut, and the order for applying bounds}\label{subsec:scale:folded}

At one cut, the comparison theorem bounds a signed difference of
records $K_s$ through the inequality \eqref{eq:record:sharp}
below: its pairing with half-functions on the two sides is at most
a constant $d$ times the product of their norms in the half
marginals. (This displayed form is what ``sharp'' means throughout
the paper.) Cauchy--Schwarz across the cut then replaces the
signed difference by the geometric mean of two nonnegative
energies, obtained by copying one half onto all slots---the folds
$\phi_\pm$ of \cref{def:sorting}. The next lemma states this
reduction, the order in which such steps may be composed, and the
marginals and prices of the two energies.

\begin{lemma}[Folded energies]\label{lem:record:folded-energy}
Let $\vartheta$ be a cut of $V$ with halves $V_\pm$ and folds $\phi_\pm$
(\cref{def:sorting}), at which every $\mu_\Blk$ is reflection positive. Let
$E\subseteq\mathcal A_N$, $\mathcal S=\mathcal A_N\setminus E$, and write
$z=(z_+,z_-)$ with $z_\pm\in X_E^{V_\pm}$. Let $\pi$ be a probability measure
on a finite set $\Omega_0$ of shared labels; for $s\in\Omega_0$ let
$\kappa_s$ be a probability measure on $X_E^V$ with half marginals $m_\pm^s$
on $X_E^{V_\pm}$, and let $K_s$ be a signed measure on $X_E^V$ such that,
for all real functions $\varphi$ on $X_E^{V_+}$ and $\psi$ on $X_E^{V_-}$,
\begin{equation}\label{eq:record:sharp}
 \Big|\int\varphi(z_+)\psi(z_-)\,dK_s(z)\Big|\le d\,
 \|\varphi\|_{L^2(m_+^s)}\|\psi\|_{L^2(m_-^s)},
\end{equation}
with $d$ independent of $s$. Let $g_v$ $(v\in V)$ be real functions on
$X_E\times X_{\mathcal S}$. Put
\begin{align*}
 T&=\int\pi(ds)\int\Lambda_{\mu_{\mathcal S}}\big((g_v(z_v,\cdot))_{v\in V}\big)\,dK_s(z),\\
 \mathfrak E_\pm&=\int\pi(ds)\int\Lambda_{\mu_{\mathcal S}}
 \big((g_{\phi_\pm(u)}(z_{\phi_\pm(u)},\cdot))_{u\in V}\big)\,dm_\pm^s(z_\pm).
\end{align*}
\begin{enumerate}[label=(\alph*)]
\item $\mathfrak E_\pm\ge0$ and $|T|\le d\,\mathfrak E_+^{1/2}\mathfrak E_-^{1/2}$.
\item $\mathfrak E_\pm$ is the energy of the inputs $g_{\phi_\pm(u)}$ with the
 nonnegative selected measure $\kappa^\pm$, the image of
 $\pi(ds)m_\pm^s(dz_\pm)$ under $z_\pm\mapsto(z_{\phi_\pm(u)})_{u\in V}$. The
 one-slot marginal of $\kappa^\pm$ at $u$ is the marginal at $\phi_\pm(u)$
 of $\kappa=\int\pi(ds)\kappa_s$. In particular, if $\kappa$ has uniform
 one-slot marginals on $V_+$ (resp.\ $V_-$), so has $\kappa^+$
 (resp.\ $\kappa^-$) on all of $V$.
\item If the inputs also depend on an external index $\omega$ with a
 probability law $\mathbf P$ independent of the labels, then
 $\E_{\mathbf P}|T^\omega|\le d\,(\E_{\mathbf P}\mathfrak E_+^\omega)^{1/2}
 (\E_{\mathbf P}\mathfrak E_-^\omega)^{1/2}$.
\item If $K_s=K'_s\otimes\pi'_s$, where $\pi'_s$ is a probability measure on
 the labels of further atoms (a coupling of its two half marginals) and
 $K'_s$ satisfies \eqref{eq:record:sharp} relative to $\kappa'_s$, then $K_s$
 satisfies \eqref{eq:record:sharp} with the same $d$ relative to
 $\kappa'_s\otimes\pi'_s$.
\item[(d$'$)] \emph{(Common cut.)} If $K_s=K^1_s\otimes K^2_s$, where $K^1_s$
 and $K^2_s$ are signed measures on the labels of two disjoint sets of atoms,
 each $K^i_s$ satisfying \eqref{eq:record:sharp} at the same cut $\vartheta$,
 relative to a probability measure $\kappa^i_s$, with the constant $d_i$,
 then $K_s$ satisfies \eqref{eq:record:sharp} relative to
 $\kappa^1_s\otimes\kappa^2_s$ with the constant $d_1d_2$; if the two factors
 have their own laws of shared labels, take the product law. By induction,
 the same holds for finitely many factors, with the product of their
 constants.
\item \emph{(Copying.)} For every function $b:V\to\R$,
 \[
  \sum_{u\in V}b(\phi_+(u))+\sum_{u\in V}b(\phi_-(u))=2\sum_{v\in V}b(v).
 \]
 Hence, for structured inputs, the slot data of $\mathfrak E_\pm$ at $u$ are
 those of the slot $\phi_\pm(u)$ of $T$, and
 $\Pi(\mathfrak E_+)\Pi(\mathfrak E_-)=\Pi^2$, where $\Pi$ is the coarse price
 of the slot data of $T$; likewise for the main prices and for every product
 $\prod_vc_v$ of nonnegative numbers attached to the slots. If the slot data
 depend on further data, such as the conductors of an index support, that
 are held fixed through a finite sequence of such steps, each applied to an
 energy produced by the previous ones, then for each fixed value of these
 data the final energies $\mathfrak E_i$, which enter the resulting bound
 with the exponents $2^{-j_i}$ ($j_i$ the number of steps leading to
 $\mathfrak E_i$), satisfy $\prod_i\Pi(\mathfrak E_i)^{2^{-j_i}}=\Pi$.
\end{enumerate}
\end{lemma}

\begin{proof}
(a) A sharp bound is a block-Gram bracket. To see this, fix $s$, identify
minus-half labels with plus-half labels through the reflection, and let
$D_\pm^s$ be the diagonal matrices of the half marginals $m_\pm^s$ and
$C_s=C_\vartheta(K_s)$. Inequality \eqref{eq:record:sharp} is equivalent to
\[
 \begin{pmatrix}dD_+^s&C_s\\C_s^{\mathsf T}&dD_-^s\end{pmatrix}\succeq0:
\]
the sharp bound makes its quadratic form nonnegative by $2xy\le x^2+y^2$;
conversely, applying a nonnegative quadratic form to $(t\varphi,\psi)$ and
varying $t$ gives the sharp bound. This also covers zero marginal masses
and $d=0$.

Let $Q=C_\vartheta(\mu_{\mathcal S})\succeq0$, by reflection positivity and
\eqref{eq:cc:kron}. The bracket with every block equal to $Q$ is positive
semidefinite. Tensorizing it with the preceding bracket
(\cref{lem:record:block-gram}(b)) gives
\[
 \begin{pmatrix}
 dD_+^s\otimes Q&C_s\otimes Q\\
 C_s^{\mathsf T}\otimes Q&dD_-^s\otimes Q
 \end{pmatrix}\succeq0.
\]
Evaluate at the two half-functions carrying their respective inputs $g_v$.
By \eqref{eq:cc:bilinear}, the off-diagonal form is the integrand defining
$T$ at fixed $s$; the diagonal forms, without the factor $d$, are exactly
the two folded energies at fixed $s$, since a diagonal copy repeats the
same half assignment across the cut. They are nonnegative, and the
block-Gram inequality bounds the absolute value of the first form by $d$
times their geometric mean. Integrating over $s$ and applying
Cauchy--Schwarz in $L^2(\pi)$ proves (a).

(b) The first statement is the definition; for the marginal, the law of
$z_{\phi_+(u)}$ under $\pi(ds)m_+^s(dz_+)$ is the marginal at
$\phi_+(u)\in V_+$ of $\int\pi(ds)\kappa_s$, because $m_+^s$ is the
$V_+$-marginal of $\kappa_s$.

(c) Apply (a) for each $\omega$ and Cauchy--Schwarz in $L^2(\mathbf P)$.

(d$'$) Tensorize the two sharp brackets from (a). Their diagonal blocks
are $d_1d_2D^1_\pm\otimes D^2_\pm$ and their off-diagonal block is the cut
matrix of $K_s^1\otimes K_s^2$, by \eqref{eq:cc:kron}. The same equivalence
therefore gives the sharp bound with constant $d_1d_2$. The argument is
pointwise in all shared labels, so independent shared-label laws can be
combined by taking their product, and induction gives finitely many factors.
For (d), a probability coupling satisfies the sharp bound with constant one
by Cauchy--Schwarz under that coupling; apply (d$'$).

(e) Each $v\in V_+$ occurs exactly twice among $\phi_+(u)$, $u\in V$ (for
$u=v$ and $u=\vartheta v$), and no $v\in V_-$ occurs; symmetrically for
$\phi_-$. This is the identity, and the same count gives
$\prod_uc_{\phi_+(u)}\cdot\prod_uc_{\phi_-(u)}=\prod_{v\in V_+}c_v^2
\prod_{v\in V_-}c_v^2=(\prod_vc_v)^2$ for all numbers $c_v\ge0$. With $c_v$
the slot factors $\xi_v^{a/72}$, $\kappa_v^{\theta/72}$ and
$(\bar\alpha_v/\alpha)^{\theta/72}$ (the last two at changed slots), this is
the statement on prices; the changed or unchanged status of a slot is copied
with its input. For a sequence of
steps, induct on the number of steps: replacing an energy $\mathfrak E$ of
exponent $2^{-j}$ by the two energies of a further step, each of exponent
$2^{-j-1}$, does not change the product, because
$\Pi(\mathfrak E_+)^{1/2}\Pi(\mathfrak E_-)^{1/2}=\Pi(\mathfrak E)$ by the
first part.
\end{proof}

By (a), (e) and \cref{prop:record:density-structured}, a sharp difference
with constant $d$ costs at most
$d(1+u_\Lsc+(C_{\rm dc}/C_0)\Lsc^{a-\sigma}\alpha^{1-\theta})\Pi\,Y$ at a least
normalized counterexample: the geometric mean of the two energies retains the
price $\Pi$, not $\Pi^2$. The same holds for the weights carried by the slots
of an index support. In the index energies of Part~III each query conductor
occurs at two slots, so it enters with the exponent $2a/72=a/36$. A later
chamber fold can give four occurrences on one energy and none on the other,
but the geometric mean keeps the original two; the conductors are therefore
held fixed through every Cauchy--Schwarz step and averaged only after the
complementary energies are combined, and no higher conductor moment is used
(\cref{lem:envelope-energy}).

The two energy estimates apply to the energies of the comparison once the
selected laws and the bounds applied to them have the following properties.
They are properties of the construction, verified where the laws are built.

\begin{definition}[Admissible laws]\label{def:admissible}
Consider a reduction of a multilinear form to structured energies by
expansions over selected atoms, sharp operator bounds at cuts
(\cref{lem:record:folded-energy}) and \cref{prop:record:density-structured}.
Its closure obligations are:
\begin{enumerate}[label=(O\arabic*)]
\item every selected law that remains after the operator-to-energy reduction
 is nonnegative, or an invariant signed measure, with uniform whole-slot
 marginals on the selected atoms (\ref{it:record:R3},
 \cref{def:record:comparison-law}), and it is independent of the interval,
 residue and inclusion indices;
\item every sharp operator bound is uniform in the shared labels at which it
 is applied (the constant $d$ of \eqref{eq:record:sharp});
\item index-support projections are taken before the field folds, and the
 inclusion conductors stay fixed through all folds and geometric means; they
 are averaged only afterwards (\cref{lem:record:folded-energy}(e));
\item the chamber folds used in a geometric mean commute with the reserved
 cut $\iota$, so that complementary folds have the
 copying property of \cref{lem:record:folded-energy}(e).
\end{enumerate}
A signed measure on the labels of a finite set of atoms is an
\emph{admissible law} if it satisfies (O1) and every sharp operator bound
applied to it satisfies (O2). A reduction is \emph{admissible} if its selected
laws are admissible and it satisfies (O3) and (O4).
\end{definition}

For the slot array $V$ every chamber fold commutes with the reserved cut
$\iota$ (\cref{lem:slot-array}(iii)); so a law fixed samplewise by $\iota$,
and an assignment of inputs invariant under $\iota$, stay so after a chamber
fold. Whole-slot uniformity does not imply uniformity of the conditional
distributions of the non-shared labels of one half given the shared labels,
and none of the statements of this section needs it. A pointwise statement
with arbitrary pinned values of the non-shared labels would need a further
theorem and is not used. Each construction lemma of Part~V states that its
laws are admissible; \cref{tab:admissibility} collects them.

\subsection{Forced pairs and very large pinnings are negligible}\label{subsec:scale:cancellation}

The accepted queries enter through one estimate: an average in which an
ordered pair of slots carries inputs from adjacent intervals, with residues
differing by a nonvanishing square class, and whose law forces a strict-square
difference on that pair at every prime with $e_p=1$ (at a clique prime, a
nonzero difference of even valuation with a square unit part), is
negligible. Square-difference-free%
ness is used here, through the actual-square kernel.

\begin{lemma}[Accepted-pair cancellation with a density factor]
\label{lem:record:density-can}
There is a constant $C_{\rm can}$ with the following property. Let
$\Lsc\ge\Lsc_*$ (\cref{lem:scale-conditions}), $N=\Kint\Lsc$, let
$A\subseteq[N]$ be square-difference-free and $\alpha=|A|/N$. Let $X_{\rm pre}$
be the product of the groups of the atoms below $\Pfive$, and let
$\varrho=\varrho_{\rm pre}\otimes\bigotimes_{\Pfive\le p\le N^\beta}\varrho_p$
be a product of probability measures on $X_{\rm pre}^V$ and on $\Fp^V$, each
with uniform one-slot marginals. Let $(u,v)$ be an ordered pair of slots
such that:
\begin{enumerate}[label=(\alph*)]
\item for every $p\ge\Pfive$, $\varrho_p$ satisfies the hypothesis of
 \cref{lem:tr:pairsupport} at $(u,v)$ with the exponent $\gamma$: the
 $(u,v)$-marginal of $\varrho_p$ is $\Unif(\Om)$, and its conditional
 operator from the pair $(u,v)$ to the other slots fixes constants, maps
 mean-zero functions to mean-zero functions and has norm at most
 $p^{-\gamma}$ on them;
\item $\varrho_{\rm pre}$-almost surely, $z_{v,p}-z_{u,p}$ is a nonzero square
 at every prime with $e_p=1$ below $\Pfive$, and at every clique prime it is
 nonzero, of even valuation $2j<e_p$, with a square unit part.
\end{enumerate}
Let $1\le b\le\Kint-1$, $s\in\Z/\Mzero\Z$, and let $r\in\Z/\Mzero\Z$ be a
nonvanishing square class (\cref{def:nonvanishing-square-class}). Let $G_u=g^{I_b,\Mzero,s}$,
$G_v=g^{I_{b+1},\Mzero,s+r}$, and let $G_w=g^{I_w,D_{0,w},s_w}$ $(w\ne u,v)$
be inputs of \cref{def:record:functional}. Then
\[
 \Big|\int\prod_{w\in V}G_w(z_w)\,d\varrho(z)\Big|\le C_{\rm can}N^{72\delta}\alpha
 \big(N^{-\gamma t_{\rm an}}+N^{-1/120+t_{\rm an}(1/2-\gamma)}\big)
 \le2C_{\rm can}\Lsc^{-\sigma}\alpha .
\]
The bound is uniform in $b$, $s$, $r$ and the other indices, so it persists
after integration against any probability law of the indices; for
$c\varrho$ with $0\le c\le1$ it is multiplied by $c$.
\end{lemma}

\begin{proof}
\emph{Step 0.} Use the constant of \cref{prop:record:kernel20}(iii)
at the triple of \cref{cor:kernel20}, with $C_D=C_H=1$. The external
modulus and the clique moduli are fixed independently of the scale. Put $H=N^{1/40}$, $T_0=N^{t_{\rm an}}$ and $D_{\rm pre}=|X_{\rm pre}|$,
which divides $Q_{\rm pre}$. By \cref{lem:scale-conditions}(v), every prime of
$[\Pfour,\Pfive)$ is a coordinate of $G_N$; the other primes below $\Pfive$
are coordinates of $G_N$ or divide $\Mzero$.

\emph{Step 1: the pair decomposition.} Condition on the prefix labels
$x=(x_w)_{w\in V}$ of all slots, distributed by $\varrho_{\rm pre}$. Put
$\varphi_x(y_u,y_v)=G_u(x_u,y_u)G_v(x_v,y_v)$ and
$W_x=\prod_{w\ne u,v}G_w(x_w,\cdot)$, functions of the tail labels. By (a),
\cref{lem:tr:pairsupport} applies to the tail factors, so
\[
 \Big|\int\varphi_xW_x\Big|^2\le\|W_x\|_2^2\Big[T_0^{-2\gamma}\|\varphi_x\|_2^2
 +\sum_{q_U\le T_0}q_U^{-2\gamma}\|E_U\varphi_x\|_2^2\Big],
\]
with $\|\varphi_x\|_2$ for $\bigotimes_p\Unif(\Om)$, $\|W_x\|_2$ for the law of
the other slots, and $U$ running over sets of tail primes $p\ge\Pfive$.
Integrate over $x$ and apply Cauchy--Schwarz:
\[
 \Big|\int\prod_wG_w\,d\varrho\Big|\le\Big(\int\|W_x\|_2^2\Big)^{1/2}
 \Big(T_0^{-2\gamma}\int\|\varphi_x\|_2^2
 +\sum_{q_U\le T_0}q_U^{-2\gamma}\int\|E_U\varphi_x\|_2^2\Big)^{1/2}.
\]

\emph{Step 2: moments.} Under $\varrho$ every $z_w$ is uniform on $G_N$. By
H\"older with $70$ factors and \cref{lem:record:lift-moments}(c) with moments of
order $140$ (all inputs have intervals of length at least
$\Lsc\ge10N^{6\beta}$, by \cref{lem:scale-conditions}(i)),
\[
 \int\|W_x\|_2^2\,d\varrho_{\rm pre}=\int\prod_{w\ne u,v}|G_w(z_w)|^2\,d\varrho
 \le\prod_{w\ne u,v}\|G_w\|_{140}^2\le C_\epsilon N^\epsilon\alpha ,
\]
and by Cauchy--Schwarz and the fourth moments,
$\int\|\varphi_x\|_2^2\le\|G_u\|_4^2\|G_v\|_4^2\le C_\epsilon N^\epsilon\alpha$.
The density exponents are $70\cdot2/140=1$ and $2\cdot2/4=1$.

\emph{Step 3: a low pair support.} Fix $U$ with $q_U\le T_0$. The function
$E_U\varphi_x$ averages both pair variables over the tail primes outside $U$.
Fix also the pair labels $y=(y_u,y_v)$ at $U$, with
$(y_{u,p},y_{v,p})\in\Om$; let $\mathcal S'$ be the set of tail primes
$\Pfive\le p\le N^\beta$ outside $U$, and $g_1=G_u(x_u,y_u,\cdot)$,
$g_2=G_v(x_v,y_v,\cdot)$ on $X_{\mathcal S'}$. By the pair multiplier
identity of \cref{lem:square-shift}(d),
\[
 E_U\varphi_x(y)=\Psi_{\le H}+\Psi_{>H},\qquad
 \Psi_{\le H}=\sum_{q(\xi)\le H}\widehat{g_1}(-\xi)\widehat{g_2}(\xi)g(\xi),
\]
the sums running over the characters of $X_{\mathcal S'}$, with $g(\xi)$ the
unit quadratic average (\cref{def:kernel-notation}). Since every prime of
$\mathcal S'$ exceeds $64$, $|g(\xi)|\le q(\xi)^{-1/3}<H^{-1/3}$ for
$q(\xi)>H$, and Cauchy--Schwarz and Parseval give
$|\Psi_{>H}|\le H^{-1/3}\|g_1\|_2\|g_2\|_2$.

For the low part put $D=\Mzero D_{\rm pre}q_U$. By
\cref{lem:scale-conditions}(vi), $HD_{\rm pre}q_U\le Q$, so
\cref{lem:record:fibre-scale}(a), with the conditioned atoms in the role of
$E$, gives for $q(\xi)\le H$
\[
 \widehat{g_1}(\xi)=\frac D\Lsc\sum_{\substack{n\in A\cap I_b\\ n\equiv\mathfrak a_1\,(D)}}e(-n\xi),
 \qquad
 \widehat{g_2}(\xi)=\frac D\Lsc\sum_{\substack{n\in A\cap I_{b+1}\\ n\equiv\mathfrak a_2\,(D)}}e(-n\xi),
\]
where $\mathfrak a_1\equiv s$, $\mathfrak a_2\equiv s+r\pmod\Mzero$, $\mathfrak a_1\equiv x_u$,
$\mathfrak a_2\equiv x_v$ on $X_{\rm pre}$ and $\mathfrak a_1\equiv y_u$, $\mathfrak a_2\equiv y_v$ on $X_U$.
Every prime below $\Pfive$ and every prime of $U$ divides $D$; hence the
characters of $X_{\mathcal S'}$ with $q(\xi)\le H$ are exactly the elements
of $\Xi_{H,D}$ (\cref{def:kernel-notation}), since a squarefree $q\le H\le N^\beta$ coprime to $D$ has all
its primes in $\mathcal S'$. The difference $\mathfrak a_2-\mathfrak a_1$ is $r$ modulo $\Mzero$,
a nonzero square at the primes of $U$ and at the primes below $\Pfive$ with
$e_p=1$, and at the clique primes nonzero of even valuation below $e_p$ with
square unit part, by (b). With
$F_i=\one_A\one_{I_{b+i-1}}\one_{n\equiv\mathfrak a_i\,(D)}$ on $[N]$ and $\widehat F_i$
normalized by $N$, we have $\widehat{g_1}(-\xi)=(DN/\Lsc)\widehat F_1(-\xi)$ and
$\widehat{g_2}(\xi)=(DN/\Lsc)\widehat F_2(\xi)$. The supports are ordered,
since $I_b$ lies below $I_{b+1}$, and no difference is a positive square. So
\cref{prop:record:kernel20}(iii) applies in the form of \cref{cor:kernel20}, with
$D\le N^{1/20}$ (\cref{lem:scale-conditions}(vi)), $H=N^{1/40}$ and
$N$ above the thresholds of the proposition (\cref{lem:scale-conditions}(vii)):
\begin{align*}
 |\Psi_{\le H}|&=\frac{DN^2}{\Lsc^2}\Big|D\sum_{\xi\in\Xi_{H,D}}g(\xi)
 \widehat F_1(-\xi)\widehat F_2(\xi)\Big|\\
 &\le\frac{DN^2}{\Lsc^2}\,\frac{C N^{-1/120}\|F_1\|_{\ell^2}\|F_2\|_{\ell^2}}N
 =C\Kint N^{-1/120}\sqrt{\alpha_{{\rm f},1}\alpha_{{\rm f},2}},
\end{align*}
where
$\alpha_{{\rm f},i}=(D/\Lsc)\|F_i\|_{\ell^2}^2=(D/\Lsc)\#\{n\in A\cap I_{b+i-1}:n\equiv\mathfrak a_i\ (D)\}$
are the two fibre densities. The fibre modulus $D$ cancels. Since
$D\le\Lsc$ (\cref{lem:scale-conditions}(vi)),
$\alpha_{{\rm f},i}\le(D/\Lsc)(\Lsc/D+1)\le2$.

Now average over $x\sim\varrho_{\rm pre}$ and
$y\sim\bigotimes_{p\in U}\Unif(\Om)$. The labels $(x_u,y_u)$ are uniform on
$X_{\rm pre}\times X_U$, so the class $\mathfrak a_1$ is uniform among the
$D_{\rm pre}q_U$ classes modulo $D$ above $s$, and
$\E\alpha_{{\rm f},1}=\Mzero\#\{n\in A\cap I_b:n\equiv s\ (\Mzero)\}/\Lsc
\le\Mzero\Kint\alpha$. Hence
$\E(\alpha_{{\rm f},1}\alpha_{{\rm f},2})\le2\E\alpha_{{\rm f},1}\le2\Mzero\Kint\alpha$;
no independence of the two classes is used. For the high part,
Cauchy--Schwarz and Jensen's inequality ($\|g_1\|_2^4\le\E|g_1|^4$ over
$X_{\mathcal S'}$) give
$\E(\|g_1\|_2^2\|g_2\|_2^2)\le\|G_u\|_4^2\|G_v\|_4^2\le C_\epsilon N^\epsilon\alpha$,
because $(x_u,y_u)$ together with the uniform outside label is uniform on
$G_N$. Since $\|E_U\varphi_x\|_2^2=\E_y|E_U\varphi_x(y)|^2$ and
$|\Psi|^2\le2|\Psi_{\le H}|^2+2|\Psi_{>H}|^2$,
\[
 \int\|E_U\varphi_x\|_2^2\,d\varrho_{\rm pre}(x)\le4C^2\Kint^3\Mzero N^{-1/60}\alpha
 +2H^{-2/3}C_\epsilon N^\epsilon\alpha\le C'_\epsilon N^{-1/60+\epsilon}\alpha ,
\]
uniformly in $U$.

\emph{Step 4: summation.} Distinct sets $U$ have distinct products, and
$0<\gamma<1/2$, so
$\sum_{q_U\le T_0}q_U^{-2\gamma}\le1+T_0^{1-2\gamma}/(1-2\gamma)$. Inserting
Steps~2 and~3 into Step~1,
\[
 \Big|\int\prod_wG_w\,d\varrho\Big|\le C_\epsilon N^\epsilon\alpha
 \big(N^{-\gamma t_{\rm an}}+N^{-1/120+t_{\rm an}(1/2-\gamma)}\big).
\]
Take $\epsilon=72\delta$; both exponents are below $-\sigma$ by (S2) and
(S3), which proves the lemma.
\end{proof}

At a nonunit clique or external-index difference the kernel is applied in
the root-class form of \cref{prop:record:kernel20}(iii), which keeps every
pinned prime in the modulus. Its constants satisfy the same uniform bounds.
For the odd prefix with a unit external index,
\cref{prop:record:kernel20}(ii), where an even-valuation square is divided
out of the modulus, may be used in place of (iii). The fourth moment
enters with the exponent exactly $1/2$, and no strict inequality between
moment exponents is needed.

\begin{lemma}[Large products with a density factor]
\label{lem:record:density-lp}
Let $(\eta_\Blk)$ be nonnegative numbers indexed by the atoms of all scales,
with $\Pi_{\rm H}=\prod_\Blk(1+\eta_\Blk q_\Blk^{1/4})<\infty$, where $q_\Blk$
is the product of the primes of $\Blk$ with $e_p=1$. For $\epsilon>0$ there is
$C_\epsilon$ with the following property. Let $\Lsc\ge1$ with
$\Lsc\ge10N^{6\beta}$, $c\ge1$, and for every set $E\subseteq\mathcal A_N$
with $q_E>N^\tau$ let $\kappa'_E$ be a finite nonnegative measure on $G_N^V$
with $\kappa'_E(1)\le c\prod_{\Blk\in E}\eta_\Blk$ whose one-slot marginals are
multiples of the uniform law, and let $(G^E_v)_v$ be inputs
$g^{I,D_0,s}$ of \cref{def:record:functional}, possibly averaged over indices
by a probability law. Then
\[
 \sum_{E:\,q_E>N^\tau}\Big|\int\prod_vG^E_v(z_v)\,d\kappa'_E(z)\Big|
 \le c\,C_\epsilon\Pi_{\rm H}N^{-\tau/4+\epsilon}\alpha .
\]
The same bound holds if the term of $E$ is replaced by
$a_E(I_E^+I_E^-)^{1/2}$ with $0\le a_E\le c\prod_{\Blk\in E}\eta_\Blk$ and
$I_E^\pm$ energies of such inputs under probability measures with uniform
one-slot marginals. With $\epsilon=\ell_0$ the right-hand side is at most
$c\,C\,\Pi_{\rm H}\Lsc^{-\sigma}\alpha$, by (S4).
\end{lemma}

\begin{proof}
For fixed $E$ and indices, H\"older's inequality with $72$ factors in
$L^{72}(\kappa'_E)$ and the uniform marginals give
$|\int\prod_vG_v\,d\kappa'_E|\le\kappa'_E(1)\prod_v\|G_v\|_{L^{72}(G_N)}$; after
averaging over the indices, H\"older over their law gives the same product of
averaged norms. By \cref{lem:record:lift-moments}(c) with moments of order $72$,
each norm is at most $(C_\epsilon N^\epsilon\Mzero^{72}\Kint\alpha)^{1/72}$, so
the product is $C_\epsilon N^\epsilon\alpha$: the density exponents are
$72/72=1$. The same bound holds for $I_E^\pm$, hence for $(I_E^+I_E^-)^{1/2}$.
Finally, by Rankin's trick,
\[
 \sum_{E:\,q_E>N^\tau}\prod_{\Blk\in E}\eta_\Blk\le N^{-\tau/4}\sum_E
 \prod_{\Blk\in E}\eta_\Blk q_\Blk^{1/4}\le N^{-\tau/4}\Pi_{\rm H}.
\]
\end{proof}

In the comparison proof the coefficient $\eta_\Blk$ of an atom
is taken to be the sum of the prices at $\Blk$ listed in
(CP8)(c), after the operator-to-energy reduction
(\cref{lem:record:rec}(ii), \cref{def:cp8}); a spectral price entering
the sum is the price produced by this reduction, not the total variation
of the original signed law. The product
$\Pi_{\rm H}$ is then finite by the summability (c) of \cref{def:cp8}: the
atoms below $\Pfour$ are finitely many, each with a finite factor in
$\Pi_{\rm H}$, and at a tail prime $p$ the prices of the selected laws are
$s_p$, $a_p$ and the matching prices $T_{{\rm L},p}=T_{{\rm R},p}$, which
add up to at most three times $s_p+a_p\le C_{\rm tail}p^{-3/2+2a}$, so that
their sum weighted by $p^{1/4}$ converges (\cref{lem:tail-prices}(iii),
(iv)).

\subsection{From $\Theta<1$ to $s(N)\le CN^{1-a/71}$ (baseline induction)}\label{subsec:scale:induction}

It remains to convert a closing inequality at a least counterexample
into a bound for every $A$. The theorem does this by induction on
the scale: if a least counterexample would satisfy $Y\le\Theta Y$
with $\Theta(0)<1$ and only the two allowed kinds of remainder, then
no counterexample exists and $Y\le C\Lsc^{-a}\alpha^\theta$ holds in
general.

\begin{theorem}[Linear-density induction]\label{thm:record:lineardensity}
Let $0<a<\sigma$ and $0<\theta\le1$, and assume
\ref{it:record:R1}--\ref{it:record:R5}. Suppose there are $\Lsc_1\ge1$,
$z_0>0$ and a function $\Theta:[0,z_0]\to[0,\infty)$, continuous at $0$ with
$\Theta(0)<1$, such that the following recurrence hypothesis holds.
\begin{enumerate}[label=(Rec)]
\item Let $C_0\ge1$, let $(\Lsc,A)$ be a least normalized counterexample for
 $(a,\theta,C_0)$ with $\Lsc\ge\Lsc_1$, and let $z\in[0,z_0]$. If every
 structured energy that the recurrence bounds at $(\Lsc,A)$ satisfies
 $|\mathfrak E|\le(1+z)\|\rho\|\,\Pi\,\Yf(\Lsc,A)$ with its coarse price $\Pi$,
 and the sum of the absolute values of its cancellation and large-product
 remainders is at most $z\,\Yf(\Lsc,A)$, then
 $\Yf(\Lsc,A)\le\Theta(z)\,\Yf(\Lsc,A)$.
\end{enumerate}
Suppose further that these remainders consist of a number, bounded
independently of $\Lsc$ and $A$, of integrals of the form bounded in
\cref{lem:record:density-can}, integrated against probability laws of the
indices and multiplied by coefficients bounded independently of $\Lsc$ and
$A$, together with finitely many sums of the form bounded in
\cref{lem:record:density-lp}, whose number and whose constants $c$ and
$\Pi_{\rm H}$ are independent of $\Lsc$ and $A$. Then there is $C$
with
\[
 \Yf(\Lsc,A)\le C\,\Lsc^{-a}\alpha^\theta
\]
for every $\Lsc\ge1$ and every square-difference-free $A\subseteq[\Kint\Lsc]$.
Consequently $\alpha^{72-\theta}\le C\Lsc^{-a}$, and
$s(N)\le C'N^{1-a/(72-\theta)}$ for every $N\ge1$. At $\theta=1$ the
conversion is $\eps=a/71$.
\end{theorem}

\begin{proof}
Let $C_{\rm R}$ be a constant with: the remainders at any least normalized
counterexample $(\Lsc,A)$ with $\Lsc\ge\Lsc_*$ have total absolute value at
most $C_{\rm R}\Lsc^{-\sigma}\alpha$. It exists by
\cref{lem:record:density-can,lem:record:density-lp} and the hypothesis on
their number and coefficients, and it does not depend on $C_0$, because the
remainders do not involve the induction.

\emph{Choice of constants.} By continuity choose $z_1\in(0,z_0]$ with
$\Theta(z)\le(1+\Theta(0))/2<1$ for $0\le z\le z_1$. Choose
$\Lsc_0\ge\max(\Lsc_1,\Lsc_*)$ such that $u_\Lsc\le z_1/2$ and
$(C_{\rm dc}+C_{\rm R})\Lsc^{a-\sigma}\le z_1/2$ for $\Lsc\ge\Lsc_0$; this is
possible because $a<\sigma$. For the finitely many $\Lsc<\Lsc_0$, every
retained coefficient of $F_A$ has modulus at most $\alpha$, so
$|F_A|\le n_\Lsc\alpha$, with $n_\Lsc$ the number of retained characters of
$G_N$, and $\Yf(\Lsc,A)\le(n_\Lsc\alpha)^{72}\le n_\Lsc^{72}\alpha$. Put
$C_{\rm base}=\max_{\Lsc<\Lsc_0}n_\Lsc^{72}$ and
$C_0=\max(1,C_{\rm base}\Lsc_0^a)$. Then
$\Yf(\Lsc,A)\le C_{\rm base}\alpha\le C_0\Lsc^{-a}\alpha^\theta$ for
$\Lsc<\Lsc_0$, since $\alpha\le\alpha^\theta$.

\emph{The least counterexample.} Suppose the bound
$\Yf(\Lsc,A)\le C_0\Lsc^{-a}\alpha^\theta$ fails for some pair, and let $\Lsc$
be the least scale at which it fails, for some $A$. Then $(\Lsc,A)$ is a least
normalized counterexample, $\Lsc\ge\Lsc_0$, $\alpha>0$ and $Y=\Yf(\Lsc,A)>0$.
By \cref{prop:record:density-structured}, every structured energy satisfies
$|\mathfrak E|\le(1+u_\Lsc+C_{\rm dc}\Lsc^{a-\sigma}\alpha^{1-\theta}/C_0)
\|\rho\|\Pi Y\le(1+z_1)\|\rho\|\Pi Y$, because $\alpha\le1$ and $C_0\ge1$. The
remainders are at most
$C_{\rm R}\Lsc^{-\sigma}\alpha=(C_{\rm R}/C_0)\Lsc^{a-\sigma}\alpha^{1-\theta}
\cdot C_0\Lsc^{-a}\alpha^\theta<z_1Y$. By (Rec), $Y\le\Theta(z_1)Y<Y$, a
contradiction. Hence the bound holds with $C=C_0$. No lower bound on $\alpha$
is used: the relative error $C\Lsc^{a-\sigma}\alpha^{1-\theta}/C_0$ is uniform
in $\alpha$ because $\theta\le1$.

\emph{Conversion.} By \cref{lem:record:holder}(e), $\alpha^{72}\le\Yf(\Lsc,A)
\le C_0\Lsc^{-a}\alpha^\theta$, so $\alpha^{72-\theta}\le C_0\Lsc^{-a}$ when
$\alpha>0$, and
$|A|=\alpha N\le C_0^{1/(72-\theta)}\Kint^{a/(72-\theta)}N^{1-a/(72-\theta)}$.
For an arbitrary $N\ge\Kint$ put $N'=\Kint\lceil N/\Kint\rceil<2N$; a
square-difference-free subset of $[N]$ is one of $[N']$, which gives
$s(N)\le C'N^{1-a/(72-\theta)}$. For $N<\Kint$, $s(N)\le N\le\Kint$.
\end{proof}

The hypothesis (Rec) is supplied by the comparison theorem: in the proof of
\cref{thm:comparison}, \cref{lem:record:rec} shows that (Rec) holds for the
density exponent $\theta$ of the frame data, with $\Theta(z)$ the closing
number of \cref{thm:comparison} evaluated with every energy bound inflated by
the factor $1+z$ and every remainder bounded by $zY$; it also proves the
continuity at $z=0$ and that the remainders have the forms of
\cref{lem:record:density-can,lem:record:density-lp}. In the baseline theorem
$\theta=1$, and $\Theta(0)$ is the certified closing value. The induction uses
only the moments of orders $2$, $4$, $72$, $140$ and $142$, and no limit in an
interpolation parameter.

\begin{remark}[Fixed parameters]\label{rem:fixed-parameters}
Every statement and proof of \cref{sec:scale-comparison} holds for an
arbitrary fixed integer $\Kint\ge2$, with the least square step
\eqref{eq:square-step}, and for arbitrary fixed even exponents $e_p=2m_p$ at
the clique primes, and so do the statements of Part~II that it uses:
the functional is defined for these data (\cref{def:record:functional}); the
index averages hold for every $\Kint\ge2$ (\cref{lem:twointerval}); the
moment constants depend on the prefix only through $Q_{\rm cl}$
(\cref{lem:record:lift-moments}(c)); conductor decay allows any finite
prefix (\cref{lem:record:prefix-conductor}), and the pair-support decay
involves only the tail primes (\cref{lem:tr:pairsupport}); and the kernel
allows any $m_p$, with constants depending on $\prod_{p\in\mathcal P_1}p^{e_p}$
(\cref{prop:record:kernel20}(iii), \cref{lem:kernel-depth-classes}).
Constants and the starting scale depend on these fixed data, and nothing is
asserted uniformly as $\Kint$ or the prefix grow with $N$. Each interval pair
$I_b,I_{b+1}$ with $1\le b<\Kint$ is ordered, which, with their common length
$\Lsc$, is all that the cancellation uses of the intervals, and the child
scales are strictly smaller by
\cref{lem:record:fibre-scale}(d). The baseline theorem uses $\Kint=\PushK$,
$e_p=2\mcl$ with $\mcl=8192$ at the class-one primes below $\Pfour$,
$\theta=1$ and the price exponent $a=\PushA$, which satisfies the gate
$a<\sigma$ of \cref{thm:record:lineardensity}, with
$\sigma-a=65195254/201416015625$, and gives
$\eps=a/71=\PushExponent$. The gates of the analytic transfer that involve
$a$ (the analytic, density, fixed-index and fixed-prefix gates,
\cref{cert:record:analytic-stage}(S0)) hold on the
whole range $0<a\le9/10000$ and for every fixed $\Kint\ge2$; in particular,
there the induction slack is $\sigma-a\ge103/330000$, and the exponent of the
bound $2^{513-\mcl+1024\mcl a}$ for the price of the deleted clique exception
(\cref{lem:frontier:iidenergy}) is at most $-80783/625$, so that this price
together with the tail prices beyond $\Pfour$ stays below $10^{-12}$
(\cref{cert:record:frame-data}(a), (f); \cref{rem:prefix-digits}).
\end{remark}

\begin{remark}[Empty and concentrated fibres]\label{rem:record:fibres}
\leavevmode
\begin{enumerate}[label=(\roman*)]
\item If $A=\emptyset$, every lift and every energy vanishes; the induction
 treats this case before any division by $\alpha$.
\item If $A$ is a single point, $\E_I|f|^2=1/|I|$, and
 \cref{lem:record:lift-moments} bounds every fixed moment by
 $C_\epsilon N^\epsilon/|I|$; no hypothesis $\alpha|I|\to\infty$ is used.
\item A pinned lift need not vanish on an empty fibre. Let $|I|$ be divisible
 by $15$, $A=\{n\in I:15\mid n\}$, $\mathcal P=\{3,5\}$ and $Q=5$. The lift is
 $G(z_3,z_5)=(3\cdot\one_{z_3=0}+5\cdot\one_{z_5=0}-1)/15$. At $z_3=1$ the
 fibre is empty, but $G(1,z_5)=(5\cdot\one_{z_5=0}-1)/15\ne0$. Its part of
 conductor at most $Q/3$ is zero, as the exact fibre identity
 (\cref{lem:record:fibre-scale}(a)) requires. The comparison never asserts
 more than this: the square-step model $\widetilde g$ vanishes on an empty
 fibre, and the difference $g-\widetilde g$ is controlled by averaged
 moments.
\item If $A$ is one residue class modulo a retained conductor, the lift
 concentrates on one class of the product group, and its moments of every
 fixed order are of order $\alpha$; so the linear dependence on $\alpha$ in
 \cref{lem:record:lift-moments} cannot be improved in general.
\end{enumerate}
\end{remark}

\begin{remark}[The density exponent]\label{rem:record:theta-ladder}
At $\theta=1$ the density saving $a/(72-\theta)$ is $a/71$. Without the
density-linear moments of \cref{lem:record:lift-moments}, a remainder
$C\Lsc^{-\sigma}$ is compared with $Y>C_0\Lsc^{-a}\alpha^\theta$ through the
sparse-set bound $Y\le C_\epsilon\Lsc^\epsilon\alpha^s$ (every $\epsilon>0$),
$s=(10^6-36)/(10^6-1)$, which follows from H\"older, Parseval and a moment of
order $2\cdot10^6$ by interpolation. This needs the sparse gate $as/(s-\theta)<\sigma$, which allows
exactly $\theta<s(\sigma-a)/\sigma$: about $0.2680$ at the price exponent of
the baseline theorem, so a saving of less than $a/71.73$. That route also needs
the uniform whole-slot marginals: its main-term Jensen step uses them
(\cref{def:admissible}). Only $\theta=0$, with
saving $a/72$, needs no density averaging at all. The density-linear moments
thus gain the factor $72/71$ over $\theta=0$, and about $1.010$ over the
sparse gate.
\end{remark}

\paragraph{Consequences and what to read next.}
Inputs used in this section: the functional, its atoms, inputs and
reference hypotheses
(\cref{def:record:functional,hyp:record:reference}); the least square step
and the strict square classes (\eqref{eq:sq}, \cref{def:strict-square-class,def:nonvanishing-square-class});
the products of atom
laws with the signed H\"older inequality, the seminorm, invariance and
averaging (\cref{lem:record:holder}); the density-linear moment bound
(\cref{lem:record:lift-moments}); conductor decay with a finite prefix and
the pair-support decay (\cref{lem:record:prefix-conductor,lem:tr:pairsupport},
with \cref{def:mixing}); the pair multiplier of the strict-square pair law
(\cref{lem:square-shift}(d)); the actual-square kernel and its notation
(\cref{def:kernel-notation,prop:record:kernel20,cor:kernel20}); the tail prices, for the
finiteness of $\Pi_{\rm H}$ (\cref{lem:tail-prices}); reflection
positivity, folds and the one-cut
Cauchy--Schwarz inequality (\cref{def:rp,def:sorting,lem:holder:one-cut}); the
slot array (\cref{lem:slot-array}); and the stage inequalities of
\cref{cert:record:analytic-stage}.
